%------------------------------------------------------------------------------------
% ВВЕДЕНИЕ
%------------------------------------------------------------------------------------
\newpage
\begin{center}
  \textbf{\large ВВЕДЕНИЕ}
\end{center}
\addcontentsline{toc}{chapter}{ВВЕДЕНИЕ}

В настоящий момент актуальной задачей развития телекоммуникаций является создание и развитие объединенных сетей
наземной сотовой связи (TN, Terrestrial Networks) и неназемной сотовой связи (NTN, Non-Terrestrial Networks).
Наземная сотовая связь состоит из вышек базовых станций, в то время как неназемная связь базируется в первую очередь
на спутниках низкой околоземной орбиты (LEO, Low Earth Orbit). В дальнейшем планируется использование высотных платформ (HAPS, High Altitude Platform Stations)
— аэростатах, дирижаблях и беспилотных летательных аппаратов (UAVs, Unmanned Aerial Vehicles) совместно со спутниками [\ref{Figaro2026}].
Средства NTN позволяют разрешить многие задачи, с которыми столкнулись при разработке наземных систем связи.
Например, обеспечение полного глобального покрытия связью всей поверхности Земли,
обеспечение более качественной передачи данных на конкретных участках (городская застройка, транспортные коридоры),
обслуживание потребностей интернета вещей (IoT), цифровых двойников и систем мониторинга [\ref{Rinaldi2020}].

NTN будет особенно полезен в условиях плотной городской застройки (для разгрузки наземных вышек),
в сельской и удаленной местности (например, Камчатка, Сибирь, Дальний Восток),
а также для обеспечения связью подвижных объектов: пассажирских самолетов, морских судов и поездов.
Внедрение спутникового сегмента в архитектуру сетей 5G позволяет решить проблему цифрового неравенства,
обеспечивая широкополосным доступом территории с низкой плотностью населения,
где развертывание наземной инфраструктуры экономически нецелесообраз\-но. 
Кроме того, космическая связь выступает в качестве физически распределенного резервного канала, 
устойчивого к наземным техногенным авариям и стихийным бедствиям, что критически повышает отказоустойчи\-вость 
национальной телекоммуникационной инфраструктуры, особенно в условиях чрезвычайных ситуаций.

Современные системы спутниковой мобильной связи базируются на технологии стандарта 3GPP 5G NR (New Radio).
На основании этого стандарта был разработан специализированный стандарт спутниковой связи 5G NR NTN [\ref{TR38811}].
Сообщество 3GPP вело активные разработки в этом направлении, и на данный момент (вплоть до Release 20)
определены ключевые архитектурные решения. В Release 17 была заложена основа для NTN 
с ретрансляционной полезной нагрузкой спутников и поддержка сценариев IoT-NTN.
В Release 18 (5G-Advanced) были внесены улучшения для эфемерид, компенсации доплеровского сдвига на абонентском оборудовании,
поддержка работы в Ka-диапазоне и повышение мобильности/непрерывности обслуживания.
Release 19 вводит регенеративную полезную нагрузку (базовая станция на борту), а
Release 20 закладывает основы для 6G-NTN.

В России создание национальных спутниковых сетей является критически важной задачей, которой
занимаются такие компании как ГК «Роскосмос»,
АО «Решетнёв» и «Бюро 1440». АО «Решетнёв» разрабатывает систему «Марафон IoT»
(132 спутника для передачи данных, связанных с интернетом вещей).
«Бюро 1440» реализует проект низкоорбитальной группировки «Рассвет» с поддержкой 5G NTN.
ГК «Роскосмос» координирует создание среднеорбитальной системы «Скиф» для широкополосного доступа в интернет.
Развитие этих проектов необходимо для технологического суверенитета РФ в области спутниковой связи.

Современные технологии космической связи требуют высокой надежности и эффективности передачи данных,
что делает актуальным исследование различных факторов, влияющих на качество сигналов.
Одним из таких факторов является эффект Доплера, который возникает в результате относительного
движения источника и приемника сигналов. В системах связи с ортогональным частотным разделением каналов [\ref{steputin20216g_v1}]
(OFDM, Orthogonal Frequency Division Multiplexing), влияние эффекта Доплера может существенно сказаться
на характеристиках принимаемого сигнала, таких как уровень искажений, скорость передачи данных и устойчивость к помехам.

Модуляция OFDM является одной из наиболее перспективных технологий для реализации высокоскоростной передачи
информации в условиях ограниченной полосы пропускания и каналов связи с замираниями.
Благодаря своей способности эффективно противостоять межсимвольной интерференции (ISI, Inter-Symbol Interference)
и многолучевому распространению, OFDM широко применяется в системах связи, включая стандарт 5G NR.
Однако при наличии значительных скоростей движения приемников или передатчиков,
характерных для спутников на низкой околоземной орбите,
эффект Доплера приводит не только к смещению несущей частоты,
но и к масштабированию спектра. Это явление вызывает потерю ортогональности между поднесущими,
что в свою очередь порождает межканальную интерференцию (ICI, Inter-Carrier Interference).

Наряду со смещением частот, для борьбы с которым существует множество классических методов
(пилот-сигналы, автоподстройка частоты), эффект масштабирования спектра является более сложной проблемой.
Так как основной принцип OFDM заключается в строгой ортогональности поднесущих, данный метод модуляции 
особенно чувствителен к ICI, которая приводит резкому снижению помехоустойчи\-вости при росте отношения сигнал/шум.

Многие исследования были направлены на изучение влияния доплеровского сдвига и масштабирования спектра
на производительность систем связи, а также на разработку методов компенсации данных эффектов [\ref{Figaro2023}, \ref{Asciolla2025}, \ref{Gannon2022}].
В научной литературе предлагаются различные подходы: от использования временного окна
и частотной коррекции до использования частичного БПФ и адаптивной фильтрации.
Однако большинство работ либо сфокусированы на наземных сценариях высокоскоростного транспорта (поезда до $500$ км/ч),
либо рассматривают спутниковые каналы в общем виде, без привязки к жестким спецификациям
физического уровня 5G NR NTN. Поэтому важно провести исследование влияния эффекта доплеровского
масштабирования спектра на реальную систему связи, на примере стандарта 5G NR-NTN.

\textbf{Гипотеза исследования.} Предполагается, что существующие в спецификации 5G NR NTN наборы параметров,
а именно количество поднесущих, по-разному чувствительны к эффекту масштабирования.
Существуют такие условия, которые приводят к существенному снижению эффективности системы.

%------------------------------------------------------------------------------------
% ПОСТАНОВКА ЗАДАЧИ
%------------------------------------------------------------------------------------
\newpage
\begin{center}
  \textbf{\large ПОСТАНОВКА ЗАДАЧИ}
\end{center}
\addcontentsline{toc}{chapter}{ПОСТАНОВКА ЗАДАЧИ}

\textbf{Целью} работы является исследование влияния эффекта масштабиро\-вания спектра на 
производительность системы спустниковой связи на примере системы стандарта 5G NR NTN.

\textbf{Задачи:}
\begin{enumerate}
    \item Провести обзор литературы;
    \item получить теоретическию оценку ICI;
    \item провести численное моделирование системы;
    \item оценить влияние эффекта масштабирования спектра на производитель\-ность системы.
\end{enumerate}

